% from ega4-2-fr.tex
\oldpage[IV-2]{206}
\item[\emph{(ii)}] (aussi notée $(\mathrm{ii}_n)$) $C$ vérifie $(\mathrm S_n)$.
\item[\emph{(iii)}] $C$ est un anneau de Cohen--Macaulay.
\item[\emph{(iv)}] $C$ est un anneau régulier.
\item[\emph{(v)}] (aussi notée $(\mathrm v_n)$) $C$ est intègre, intégralement
clos et vérifie $(\mathrm R_n)$.
\item[\emph{(vi)}] $C$ est intègre et intégralement clos.
\item[\emph{(vii)}] $C$ est intègre.
\item[\emph{(viii)}] $C$ est réduit.
\end{enumerate}
Toutes ces propriétés vérifient $(\mathrm R_{\mathrm{III}})$, comme on l'a vu
dans (7.3.19, (i)). Elles vérifient aussi $(\mathrm R_{\mathrm{IV}})$ : pour
(i), cela résulte de (0, 16.4.10, (ii)), et pour (ii) et (iii), cela résulte de
(5.12.4) et du fait qu'un anneau local noethérien complet est caténaire
(5.6.4); pour (iv), c'est un cas particulier de (0, 17.1.8); pour (vii), cela
découle de (3.4.5) et pour (viii) de (3.4.6); pour (vi) cela résulte de
(5.12.7) et de ce qu'un anneau local noethérien complet est caténaire (5.6.4);
enfin, pour (v), cela résulte de ce qui précède et de (5.12.5).

On peut donc appliquer le théorème (7.5.1) lorsque $\boldsymbol R$ est l'une
quelconque des propriétés ci-dessus. D'autre part, on a déjà remarqué (7.3.11)
que les propriétés (i) à (viii) vérifient $(\mathrm R_{\mathrm{II}})$ (sauf
pour (vii), où la vérification de $(\mathrm R_{\mathrm{II}})$ est conséquence
de (2.1.14)). En ce qui concerne $(\mathrm R_{\mathrm V})$, on notera que si
l'on y fait $\boldsymbol R'=\boldsymbol R$, la condition
$(\mathrm R_{\mathrm V})$ se réduit à la condition
$(\mathrm R'_{\mathrm I})$ de (7.3.11), et on a noté que les propriétés (ii) à
(vi), ainsi que (viii), vérifient $(\mathrm R'_{\mathrm I})$ (7.3.11); pour la
propriété (i), on peut prendre pour $\boldsymbol R'$ la propriété d'être un
anneau de Cohen--Macaulay, en vertu de (6.3.2). Enfin, pour la propriété (vii),
la condition $(\mathrm R'_{\mathrm I})$ n'est plus vérifiée (ni d'ailleurs
$(\mathrm R_{\mathrm I})$), comme le montre l'exemple de (6.15.11, (ii)) ou de
(6.5.5, (ii)). Par contre, $(\mathrm R_{\mathrm V})$ est alors vraie lorsque
l'on prend pour $\boldsymbol R'$ la propriété d'être régulier : il suffit en
effet d'appliquer le lemme (7.5.1.1) en prenant pour $\mathcal C$ la catégorie
de tous les anneaux locaux noethériens, et pour $\boldsymbol R(C)$ la propriété
d'être intègre; cela est possible en vertu de (3.4.5).

On voit ainsi en particulier que lorsque $\boldsymbol R$ est l'une quelconque
des propriétés (i) à (viii), la conclusion du corollaire (7.5.2) est applicable
lorsque l'on suppose que les fibres formelles de $A$ sont géométriquement
régulières (cf. (7.8.2)).
\end{examples}

\begin{remarks}[7.5.4]
\EGAReaderTarget{ega4.c04.item.7-5-4}{IV.2.s7.5.4-fr}
\begin{enumerate}
\item[\emph{(i)}] Il serait intéressant de savoir si le théorème (7.5.1)
subsiste sans l'hypothèse de finitude (i) sur le corps résiduel de $B$. La
réponse est affirmative lorsque le corps résiduel $k$ de $A$ est de
caractéristique 0, comme on le voit en utilisant les résultats de Hironaka sur
la résolution des singularités (cf. (7.9.8)). Signalons le cas particulier
suivant de la question soulevée ici : Soient $A$ et $B$ deux anneaux
noethériens complets, $k$ le corps résiduel de $A$,
$\varphi:A\to B$ un homomorphisme local faisant de $B$ une $A$-algèbre
formellement lisse (0, 19.3.1) (ce qui équivaut à dire que $B$ est un
$A$-module plat et que $B\otimes_A k$ est géométriquement régulier sur $k$
(0, 19.7.1)). Alors les fibres du morphisme
${}^a\varphi:\operatorname{Spec}(B)\to\operatorname{Spec}(A)$ sont-elles
géométriquement régulières ? Il en est ainsi lorsque le corps résiduel $k'$ de
$B$ est une extension finie de $k$, d'après (7.5.1); on peut démontrer que la
réponse est encore affirmative lorsque $k'$ est une extension de type fini de
$k$. Nous ignorons la réponse lorsque $B\otimes_A k$ est un corps égal à la
clôture séparable de $k$.

\item[\emph{(ii)}] On pourrait énoncer un résultat analogue à (7.5.1) relatif
à un $A$-module $M$, un $B$-module $N$, tous deux de type fini, concernant les
propriétés de
\[
(M\otimes_A N)_y\otimes_{\mathscr O_x}k(y),
\]
où $y$ parcourt $\operatorname{Spec}(B)$ et $x=f(y)$, découlant de propriétés
de même type de $M_x$ et $N_y$.

\item[\emph{(iii)}] Lorsque la propriété $\boldsymbol R$ vérifie les conditions
$(\mathrm R'_{\mathrm I})$, $(\mathrm R_{\mathrm{II}})$,
$(\mathrm R_{\mathrm{III}})$ et $(\mathrm R_{\mathrm{IV}})$ et que les
hypothèses (i) et (ii) de (7.5.2) sont remplies, alors, des propriétés
$\boldsymbol R(A)$ et
\[
\boldsymbol P(\operatorname{Spec}(\widehat B
\otimes_{\widehat A}k),k)
\]
on déduit $\boldsymbol R(B)$. C'est le cas, comme on l'a vu (7.5.3), pour les
propriétés données en exemples dans (7.5.3), sauf pour (i) et (vii). En ce qui
concerne (vii), il est cependant plausible que la réponse à la question
suivante soit affirmative : soit $\varphi:A\to B$ un homomorphisme local
d'anneaux locaux noethériens, faisant de $B$ un $A$-module plat. Supposons que
$A$ soit complet, intègre et géométriquement unibranche et que la fibre
$\operatorname{Spec}(B\otimes_A k)$ (où $k$ est le corps résiduel de $A$) soit
géométriquement localement intègre; alors est-il vrai que $B$ soit intègre ?
On peut le démontrer lorsque $k$ est de caractéristique 0, en utilisant la
résolution des singularités de Hironaka (7.9.8); on peut montrer aussi que la
réponse est affirmative lorsque $B$ est une $A$-algèbre essentiellement de
type fini (1.3.8) (cf. (11.3.10) et (11.3.11)) ou lorsque
$\operatorname{Spec}(B\otimes_A k)$ est géométriquement normale (car il résulte
alors de (7.5.3) que le morphisme
$\operatorname{Spec}(B)\to\operatorname{Spec}(A)$ est normal, donc on conclut
par (6.15.10)). Mais même en supposant que le corps résiduel de $B$ soit égal
à celui de $A$ et que $B$ est aussi complet, nous ignorons si la réponse est
affirmative dans le cas général.

\item[\emph{(iv)}] Considérons la propriété $\boldsymbol R(C)$ : « $C$ est
réduit, équidimensionnel et vérifie $(\mathrm R_n)$ »; elle vérifie
$(\mathrm R_{\mathrm{IV}})$, en vertu de (5.12.5), mais on ignore par contre si
elle vérifie $(\mathrm R'_{\mathrm I})$, $(\mathrm R_{\mathrm{II}})$ ou
$(\mathrm R_{\mathrm{III}})$, les difficultés provenant de la vérification de
l'équidimensionalité.
\end{enumerate}
\end{remarks}

